% !TeX program = xelatex
% !TeX encoding = UTF-8
\documentclass{article}
\usepackage{ctex}
\usepackage{geometry}
\geometry{a4paper, left=3cm, right=3cm, top=2.5cm, bottom=2.5cm}

\usepackage{graphicx}
\usepackage{enumerate}
\usepackage{amsmath, amssymb, amsthm}
\usepackage{algorithm, algorithmic}
\floatname{algorithm}{算法}
\renewcommand{\algorithmicrequire}{\textbf{输入:}}
\renewcommand{\algorithmicensure}{\textbf{输出:}}

\usepackage{listings}
\usepackage{xcolor}
\usepackage{hyperref}
\hypersetup{colorlinks, citecolor=black, filecolor=black, linkcolor=black, urlcolor=black}
\usepackage{booktabs}
% 代码样式
\lstset{
    language=C++,
    keywordstyle=\color{blue},
    commentstyle=\color{green!50!black},
    stringstyle=\color{red},
    numbers=left,
    numberstyle=\tiny,
    basicstyle=\ttfamily\small,
    frame=single,
    breaklines=true,
    extendedchars=false
}

\title{第七次上机作业}
\author{学号：241830137\quad 姓名：朱子健}
\date{\today}

\begin{document}

\maketitle

\section{问题}
实现快速傅里叶变换 FFT，并应用于某个具体场景。如信号，图像，声音处理，多项
式乘法，偏微分方程求解 (谱方法)，或其它应用场景。

\section{算法原理与设计}
\subsection{算法原理}
\subsubsection{快速傅里叶变换}
离散Fourier变换（或逆变换）可以归结为计算
\[
c_{k} = \sum_{j=0}^{N-1}x_{j}\omega^{kj}    (k = 0, 1, \cdots,N-1)
\]
其中\[\omega = e^{-i\frac{2\pi}{N}}\]或\[\omega = e^{i\frac{2\pi}{N}}\],${x_{j}}(j = 0, 1, \cdots,N-1)$为已知复数序列. 计算全部 $c_{k}$ 共需 $N^{2}$ 个
操作. 当 $N $很大时计算量也很大, 直到上个世纪 60 年代中期 Cooley \& Tukey 提出了快
速傅里叶变换(FFT)算法, 可将计算工作量降到大约$O(NlogN)$, 才大大提高了运算
速度。
FFT算法的基本思想是尽量减少乘法的次数, 例如 $ab + ac + ad = a(b + c + d)$, 左端三次乘法而右端只用一次乘法. 实际上$e^{i\frac{2\pi}{N}j}(k, j = 0, 1, \cdots,N-1)$只有$N$个不同的值.特别是$N = 2^{P}$时只有$\frac{N}{2}$个不同的值。因此在计算$c_{k}$时, 可合并
大量同因子的项, 从而减少乘法次数.
\subsubsection{用快速傅里叶变换实现多项式乘法}
设两个多项式
\[
A(x)=a_0+a_1x+\cdots+a_{n-1}x^{n-1},\qquad 
B(x)=b_0+b_1x+\cdots+b_{m-1}x^{m-1},
\]
其乘积 \(C(x)=A(x)\cdot B(x)\) 的系数为卷积：
\[
c_k=\sum_{i=0}^{k}a_i b_{k-i},\quad k=0,1,\dots,n+m-2.
\]
直接计算的时间复杂度为 \(O(nm)\)。利用快速傅里叶变换（FFT）可降至 \(O(N\log N)\)，其中 \(N\) 为不小于 \(n+m-1\) 的 \(2\) 的幂。核心依据是\textbf{卷积定理}：
\[
\mathcal{F}\{A * B\} = \mathcal{F}\{A\} \cdot \mathcal{F}\{B\},
\]
即在点值表示下，多项式乘积对应点值逐点相乘。


\subsubsection{算法设计}
\begin{algorithm}[H]
\caption{快速傅里叶变换实现方法}
\begin{algorithmic}[1]
    \STATE Step 1 给出复数数组 \( A_0(N), A_1(N) \) 及 \( w \left( \frac{N}{2} \right) \)，确定 \( P, q = 1 \)。将已知数据 \(\{x_j\}\) 输入到数组 \( A_0(N) \) 的单元中。
    \STATE Step 2 计算 \( w^m = e^{-i\frac{2\pi}{N}m} \)（或 \( w^m = e^{i\frac{2\pi}{N}m} \)）（\( m \) 从 0 到 \( 2^{P-1}-1 \)），结果存放在 \( w\left(\frac{N}{2}\right) \) 的相应单元中。
    \STATE Step 3 \( j \) 从 0 到 \( (2^{P-q}-1) \)，\( k \) 从 0 到 \( (2^{q-1}-1) \)，按 FFT 公式(3.5.15)计算 \( A_q \)，结果放在 \( A_1(N) \) 的相应单元中。
    \STATE \[
    \begin{cases}
        A_1(j2^q + k) = A_0(j2^{q-1} + k) + A_0(j2^{q-1} + k + 2^{P-1}) \\
        A_1(j2^q + k + 2^{q-1}) = [A_0(j2^{q-1} + k) - A_0(j2^{q-1} + k + 2^{P-1})]w(j2^{q-1})
    \end{cases}
    \]
    \STATE Step 4 \( k, j \) 循环结束，\( q+1 \to q, A_1(N) \to A_0(N) \)。
    \STATE 若 \( q < P \)，则 goto Step 3；否则 goto Step 5。
    \STATE Step 5 此时 \( q = P, j = 0 \)，\( k \) 从 0 到 \( (2^{P-1}-1) \)。计算
    \STATE \[
    \begin{cases}
        A_1(k) = A_0(k) + A_0(k + 2^{P-1}) \\
        A_1(k + 2^{q-1}) = A_0(k) - A_0(k + 2^{P-1})
    \end{cases}
    \]
    \STATE 这里得到的 \( A_1(k)(k = 0, 1, \dots, 2^{P-1}) \) 即为所求 \( c_k \)。
\end{algorithmic}
\end{algorithm}

\begin{algorithm}[H]
\caption{快速傅里叶变换实现多项式乘法}
\begin{algorithmic}[1]
    \REQUIRE 多项式 $A(x)=\sum_{i=0}^{n-1} a_i x^i$，$B(x)=\sum_{j=0}^{m-1} b_j x^j$，系数为整数或实数
    \ENSURE 乘积多项式 $C(x)=A(x)B(x)=\sum_{k=0}^{n+m-2} c_k x^k$ 的系数数组 $c[0..n+m-2]$
    \STATE 令 $L = n + m - 1$，$N = 2^{\lceil \log_2 L \rceil}$  
    \STATE 构造数组 $A_{\text{pad}}[0..N-1]$，$B_{\text{pad}}[0..N-1]$，初值为0
    \FOR{$i=0$ 至 $n-1$} \STATE $A_{\text{pad}}[i] \gets a_i$ \ENDFOR
    \FOR{$j=0$ 至 $m-1$} \STATE $B_{\text{pad}}[j] \gets b_j$ \ENDFOR
    \STATE $\hat{A} \gets \text{FFT}(A_{\text{pad}})$  
    \STATE $\hat{B} \gets \text{FFT}(B_{\text{pad}})$
    \FOR{$k=0$ 至 $N-1$} \STATE $\hat{C}[k] \gets \hat{A}[k] \cdot \hat{B}[k]$ \ENDFOR
    \STATE $C_{\text{pad}} \gets \text{IFFT}(\hat{C})$  
    \FOR{$k=0$ 至 $L-1$}
        \STATE $c_k \gets \text{round}\bigl(\operatorname{Re}(C_{\text{pad}}[k])\bigr)$ 
    \ENDFOR
    \STATE \RETURN $c[0..L-1]$
\end{algorithmic}
\end{algorithm}

\section{结果分析}

\begin{table}[htbp]
\centering
\caption{FFT 多项式乘法示例}
\begin{tabular}{|c|l|l|l|}
\hline
编号 & 多项式 A(x) & 多项式 B(x) & 乘积 C(x) \\
\hline
1 & $1+2x+3x^2$ & $4+5x+6x^2$ & $4+13x+28x^2+27x^3+18x^4$ \\
\hline
2 & $1-x+x^2-x^3$ & $1+x+x^2+x^3$ & $1 - x^6$ \\
\hline
3 & $\sum_{i=0}^{9}(i+1)x^i$ & $\sum_{i=0}^{9}(i+1)x^i$ & \begin{tabular}[t]{@{}l@{}} $1+4x+10x^2+20x^3+35x^4+56x^5+84x^6$\\ $+120x^7+165x^8+220x^9+276x^{10}+324x^{11}+360x^{12}$\\ $+380x^{13}+380x^{14}+360x^{15}+324x^{16}+276x^{17}$\\ $+220x^{18}+165x^{19}+120x^{20}+84x^{21}+56x^{22}$\\ $+35x^{23}+20x^{24}+10x^{25}+4x^{26}+x^{27}$ \end{tabular} \\
\hline
4 & $2+3x+5x^2+7x^3$ & $11+13x+17x^2+19x^3$ & $22+59x+128x^2+254x^3+317x^4+337x^5+254x^6+133x^7$ \\
\hline
\end{tabular}
\end{table}

\section{附录: Python 源代码实现}
\begin{lstlisting}[language=Python]
import math, cmath
def fft(a, invert):
    n = len(a)
    if n == 1:
        return a
    # 分治：奇偶下标
    even = fft(a[0::2], invert)
    odd  = fft(a[1::2], invert)
    # 合并
    ang = 2 * math.pi / n * (-1 if invert else 1)
    w = 1
    wn = complex(math.cos(ang), math.sin(ang))
    for i in range(n // 2):
        t = w * odd[i]
        a[i] = even[i] + t
        a[i + n//2] = even[i] - t
        w *= wn
    if invert:
        for i in range(n):
            a[i] /= n
    return a

def poly_mul_fft(a, b):
    n = 1
    while n < len(a) + len(b) - 1:
        n <<= 1
    fa = a + [0] * (n - len(a))
    fb = b + [0] * (n - len(b))
    fft(fa, False)
    fft(fb, False)
    for i in range(n):
        fa[i] *= fb[i]
    fft(fa, True)
    return [int(round(fa[i].real)) for i in range(len(a) + len(b) - 1)]
\end{lstlisting}
\begin{lstlisting}
import numpy as np
def poly_mul_fft(a, b):
    total_len = len(a) + len(b) - 1
    n = 1
    while n < total_len:
        n <<= 1
    fa = np.fft.fft(a + [0] * (n - len(a)))
    fb = np.fft.fft(b + [0] * (n - len(b)))
    fc = fa * fb
    c = np.fft.ifft(fc).real
    return [int(round(x)) for x in c[:total_len]]
def main():
    print("多项式乘法（使用 FFT）")
    print("请输入第一个多项式的系数列表，例如 [1,2,3]：")
    a = eval(input().strip())
    print("请输入第二个多项式的系数列表：")
    b = eval(input().strip())
    c = poly_mul_fft(a, b)
    print("乘积系数：", c)
if __name__ == "__main__":
    main()
\end{lstlisting}
\end{document}